\def\level{Première}  % Classe à droite
\def\course{NSI}
\def\chaptername{Les dictionnaires} % Titre au centre
\def\docname{AP07-S}        % Identifiant à gauche

\pagestyle{cours_FP}
\makeheader{} % Affiche le bandeau

\begin{definition}[Les dictionnaires]{}
	
\begin{CodePythontexAlt}[Largeur=1\linewidth,Centre,PremLigne=1]{}
mon_dictionnaire = {}				# Dictionnaire vide
mon_dictionnaire = {'a': 1, 'b': 2, 'c': 3} 	# Dictionnaire avec des paires clé-valeur
print(mon_dictionnaire['a'])  # Affiche 1
print(len(mon_dictionnaire))  # Affiche 3
mon_dictionnaire['a'] = 4  # Modifie la valeur associée à la clé 'a'
		\end{CodePythontexAlt}
\end{definition}

\begin{easybox}[Itérer sur un dictionnaire]{}
\begin{minipage}{0.31\linewidth}
		\begin{CodePythontexAlt}[Largeur=1\linewidth,Centre,Lignes=false,PremLigne=1]{}
	# Itérer sur les clés	
	for cle in dico.keys():
		print(dico[cle])
		\end{CodePythontexAlt}
\end{minipage}\hfill
\begin{minipage}{0.33\linewidth}
		\begin{CodePythontexAlt}[Largeur=1\linewidth,Centre,Lignes=false,PremLigne=1]{}
	# Itérer sur les valeurs
	for val in dico.values():
		print(val)
		\end{CodePythontexAlt}
\end{minipage}\hfill
\begin{minipage}{0.36\linewidth}
		\begin{CodePythontexAlt}[Largeur=1\linewidth,Centre,Lignes=false,PremLigne=1]{}
	# Itérer sur les items
	for cle, val in dico.items():
		print(cle, val)
		\end{CodePythontexAlt}
\end{minipage}
\end{easybox}


\begin{easybox}{Création d'un dictionnaire}{}
	\begin{CodePythontexAlt}[Largeur=1\linewidth,Centre,PremLigne=1]{}
# Création d'un dictionnaire par itération
mon_dico = {}
for i in range(5):
	mon_dico[i] = i**2
print(mon_dico) # Affiche {0: 0, 1: 1, 2: 4, 3: 9, 4: 16}

# Création d'un dictionnaire par compréhension
mon_dico = {i: i**2 for i in range(5)}
print(mon_dico) # Affiche {0: 0, 1: 1, 2: 4, 3: 9, 4: 16}
\end{CodePythontexAlt}
\end{easybox}

\begin{easybox}{Appartenance d'un élément à un dictionnaire}{}

	\begin{CodePythontexAlt}[Largeur=1\linewidth,Centre,PremLigne=1]{}
mon_dico = {'a': 1, 'b': 2, 'c': 3}
print('a' in mon_dico.keys())  # Affiche True
print(4 in mon_dico.values())  # Affiche False
print(('a', 1) in mon_dico.items())  # Affiche True
\end{CodePythontexAlt}
\end{easybox}

\begin{easybox}{Suppression dans un dictionnaire}{}

	\begin{CodePythontexAlt}[Largeur=1\linewidth,Centre,PremLigne=1]{}
del mon_dico['a']  # Supprime la paire clé-valeur 'a': 1
mon_dico.clear()  # Supprime toutes les paires clé-valeur du dictionnaire
\end{CodePythontexAlt}
\end{easybox}

\begin{easybox}{Occurrences dans un dictionnaire}{}
\begin{CodePythontexAlt}[Largeur=1\linewidth,Centre,Lignes=false,PremLigne=1]{}
def occurence(sequence)->dict:
	dico={}
	for element in sequence:
		if element in dico:
			dico[element] += 1
		else:
			dico[element] = 1
	return dico
\end{CodePythontexAlt}
\end{easybox}

